\documentclass[11pt]{article}
\usepackage[margin=0.9in]{geometry}
\usepackage{amsmath,amsfonts,amssymb,amsthm}
\usepackage{enumitem}
\usepackage[]{graphicx}
\usepackage{color,subfigure}
\usepackage{multicol}
\usepackage{float}
\usepackage[all]{xypic}
\usepackage[colorlinks=true,citecolor=cyan,linkcolor=magenta]{hyperref}
\usepackage{colonequals}

\usepackage{stmaryrd}
\usepackage{fancyhdr, lastpage}
\pagestyle{fancy}
%\fancyfoot[C]{{\thepage} of \pageref{LastPage}}



\DeclareMathOperator{\mSpec}{mSpec}
\DeclareMathOperator{\Spec}{Spec}
\DeclareMathOperator{\Ass}{Ass}
\DeclareMathOperator{\Supp}{Supp}
\DeclareMathOperator{\height}{height}
\DeclareMathOperator{\Hom}{Hom}
\DeclareMathOperator{\ann}{ann}
\DeclareMathOperator{\End}{End}
\DeclareMathOperator{\coker}{coker}
%\DeclareMathOperator{\ker}{ker}
\DeclareMathOperator{\rank}{rank}
\DeclareMathOperator{\im}{im}
\DeclareMathOperator{\Ext}{Ext}
\DeclareMathOperator{\Tor}{Tor}
%\DeclareMathOperator{\dim}{dim}
\DeclareMathOperator{\HH}{H}

\DeclareMathOperator{\lcm}{lcm}

\def\ra{\rightarrow}
\newcommand{\m}{\mathfrak{m}}
\newcommand{\C}{\mathbb{C}}
\newcommand{\Q}{\mathbb{Q}}
\newcommand{\R}{\mathbb{R}}
\newcommand{\N}{\mathbb{N}}
\newcommand{\ov}[1]{\overline{#1}}


\newcommand{\Z}{\mathbb{Z}}
\DeclareMathOperator{\pdim}{pdim}
\DeclareMathOperator{\kos}{kos}
\DeclareMathOperator{\embdim}{embdim}
\DeclareMathOperator{\depth}{depth}
\DeclareMathOperator{\grade}{grade}
\DeclareMathOperator{\injdim}{injdim}


\title{}
\date{\vspace{-0.5in}}

\makeatletter
\g@addto@macro\@floatboxreset\centering
\makeatother

\theoremstyle{definition}
\newtheorem{problem}{Problem}


\begin{document}

\thispagestyle{fancy}
\pagestyle{fancy}
\rhead{UNL}
\lhead{Math 906}
\chead{Problem Set 4}

\vspace{3em}

\begin{center}
	{\LARGE Problem Set 4 \\}
	Due Wednesday, April 22, 2026
\end{center}

\

\noindent
{\bf Instructions:}
You are encouraged to work together on these problems, but each student should hand in their own final draft, written in a way that indicates their individual understanding of the solutions. Never submit something for grading that you do not completely understand. You cannot use any resources besides me, your classmates, and our course notes.


I will post the .tex code for these problems for you to use if you wish to type your homework. If you prefer not to type, please {\em write neatly}. As a matter of good proof writing style, please use complete sentences and correct grammar. You may use any result  stated or proven in class or in a homework problem, provided you reference it appropriately by either stating the result or stating its name (e.g. the definition of depth or Hilbert's Syzygy Theorem). Please do not refer to theorems by their number in the course notes, as that can change.


\vspace{2em}


\noindent
Turn in {\bf 4 problems} of your choosing. Any problem you do not turn in is now a known theorem.

\



\begin{problem}
Let $(R, \m)$ be a noetherian local ring. Show that if exists a nonzero finitely generated injective $R$-module $E$, then $\dim(R) = 0$. 
\end{problem}

\vfill


\begin{problem}
	Let $k$ be a field and $R = k \llbracket x \rrbracket$. Find the injective dimension and Bass numbers of $R$, and write down the modules in a minimal injective resolution for $R$.
\end{problem}

\vfill


\begin{problem}
	Let $(R, \m, k)$ be a noetherian ring. Show that if $R$ is injective then $R \cong E(R/\m)$.
\end{problem}

\vfill

\begin{problem}
	Let $(R, \m, k)$ be a noetherian local ring.
	
	\vspace{-0.5em}
	\begin{enumerate}[label=\alph*), itemsep=-0.1em]
		\item Show that $\ann(E(k)) = 0$.
		\item Give an example of such a ring $R$ and a nonzero injective $R$-module $I$ such that $\ann(I) \neq 0$.
	\end{enumerate}
\end{problem}

\vfill

\begin{problem}
	Let $k$ be a field.
	
	\vspace{-0.5em}
	\begin{enumerate}[label=\alph*), itemsep=-0.1em]
		\item Give an example (with proof) of a complete intersection ring $A$ that is not regular.
		\item Show that
	\[ B = k \llbracket x,y,z \rrbracket /(x^2-y^2, x^2-z^2, xy, xz,yz) \]
	is Gorenstein but not a complete intersection.
	\item Let $k$ be any field. Show that
	\[ C = k \llbracket x,y \rrbracket /(x^2, xy, y^2) \]
	is Cohen-Macaulay but not Gorenstein.
	\item Give an example (with proof) of a noetherian local ring $B$ that is not Cohen-Macaulay.
	\end{enumerate}
\end{problem}

\vfill

\begin{problem}
	Let $(R, \m, k)$ be an artinian local ring. Show that $R$ is Gorenstein if and only if all proper ideals $I$ satisfy $(0 : (0 : I)) = I$.
\end{problem}

\vfill

\begin{problem}
	Let $R$ be a noetherian local ring. The following two facts will be used together in class, but the proofs are unrelated.

\vspace{-0.5em}
\begin{enumerate}[label=\alph*), itemsep=-0.1em]
\item Show that for all proper ideals $I$, if $E$ is an injective $R$-module then $\Hom_R(R/I, E)$ is an injective module over $R/I$.
\item Let
\[ \xymatrix{
 0 \arrow[r] &	E_1 \arrow[r] & \cdots \arrow[r] & E_n \arrow[r] & 0
}\]
 be an exact sequence of $R$-modules with $n \geqslant 3$. Show that if $E_1, \ldots, E_{n-1}$ are all injective, then so is $E_n$.
\end{enumerate}
\end{problem}




\end{document}